\documentclass{article}
\usepackage{amsmath,amssymb}
\usepackage[margin=1in]{geometry}
\usepackage{enumitem}
\begin{document}

\section*{Tensor Omics: Full Toolkit Overview (May 7, 2025)}

\subsection*{A. Data Preprocessing and Representation}
\begin{itemize}[leftmargin=1.5em]
  \item \textbf{Normalization:} standard deviation scaling, optional quantile normalization, log transform (\( \log(1 + x) \)); fold-change computation for stress experiments.
  \item \textbf{Expression Vector Construction:} replicate averaging; vectors stored in raw-space and normalized-space K-D Trees.
  \item \textbf{Grid Construction and Smoothing:} automatic grid sizing; interpolation and smoothing of vector fields performed post-space construction and pre-shift calculation.
\end{itemize}

\subsection*{B. Semantic Annotation and Axes}
\begin{itemize}[leftmargin=1.5em]
  \item \textbf{Metadata:} gene families, duplication metadata, molecular functions, phenotype labels.
  \item \textbf{Anchor Axis:} user-defined semantic reference axis (e.g., time, tissue).
  \item \textbf{Parzival Axes:} lexical contrast vectors (e.g., stress, phenotype modules).
\end{itemize}

\subsection*{C. Shift Field and Geometric Observation}
\begin{itemize}[leftmargin=1.5em]
  \item \textbf{Shift Vectors:} defined as \( \vec{s} = \vec{p} - \vec{o} \); interpolated and smoothed post-estimation.
  \item \textbf{Outlier Detection:} via RDI (distance), RAI (angle), and RAS (axis-based signal).
  \item \textbf{Tissue Specificity:} angle to space diagonal as cosine similarity; versatility analysis.
  \item \textbf{RAP Projection:} Relative Axes Plane for visualizing rotation and preference; subspace RAP for phenotype-limited axes.
\end{itemize}

\subsection*{D. Duplication Outcome Analysis}
\begin{itemize}[leftmargin=1.5em]
  \item \textbf{Subfunctionalization:} detected via low RDI and RAI of in-paralog sum vector.
  \item \textbf{Neofunctionalization:} detected by RAS — newly activated axes relative to centroid.
  \item \textbf{Dosage Effects:} aligned but amplified in-paralog sum vector (low RAI, high RDI).
\end{itemize}

\subsection*{E. Subfield Geometry and Tensor Summary}
\begin{itemize}[leftmargin=1.5em]
  \item \textbf{Directional Subfield Selection:} based on metadata, DOI cosine threshold, and background exclusion.
  \item \textbf{Field Descriptors:}
    \begin{itemize}
      \item Flow vector (first eigenvector of \( C = MM^T \))
      \item Rotation vector (antisymmetric component \( C - C^T \))
      \item Frobenius norm (total signal strength)
      \item Rank (via SVD)
    \end{itemize}
\end{itemize}

\subsection*{F. Gradient Visualization: Gwion (formerly GAWD)}
\begin{itemize}[leftmargin=1.5em]
  \item \textbf{Directional Slicing:} trace field along its dominant flow vector.
  \item \textbf{Panel 1:} Line plot of standardized Frobenius norm and rank along axis.
  \item \textbf{Panel 2:} Clock plot of flow vectors orthogonal to the axis.
  \item \textbf{Panel 3:} Clock plot of rotation vectors, same projection.
\end{itemize}

\subsection*{G. Temporal and Trajectory Analysis}
\begin{itemize}[leftmargin=1.5em]
  \item \textbf{Trajectory Modes:} mean-plane projection (Mode A) vs. locally stacked planes (Mode B).
  \item \textbf{Derived Metrics:} Angular Shift Velocity (ASV), Instantaneous Acceleration (IA), vector distance to centroid.
\end{itemize}

\subsection*{H. Tensor Field Analysis — Bobito}
\begin{itemize}[leftmargin=1.5em]
  \item \textbf{Bobito 🦋 (Flechita) — Level One:}
    \begin{itemize}
      \item Cone-based expansion using vector alignment (JOE/JOV)
      \item Cutoff detection via derivative jump
    \end{itemize}
  \item \textbf{Bobito ☁️ (Nubecita) — Level Two:}
    \begin{itemize}
      \item Subfield grouping via directional coherence of flow/rotation vectors
      \item Empirical jump threshold based on sorted cosine similarity
    \end{itemize}
\end{itemize}

\subsection*{I. Semantic Glyph Learning — Gilgamesh}
\begin{itemize}[leftmargin=1.5em]
  \item \textbf{Gilgamesh Uruk (The Divider):}
    \begin{itemize}
      \item Trains local SVMs on condition-labeled subfields
      \item Tests empirical significance via permutation or bootstrap
      \item Enables glyph recognition from semantic separation
    \end{itemize}
  \item \textbf{Gilgamesh Humbaba (The Constellar):}
    \begin{itemize}
      \item Aligns subfields by Brokk axes
      \item Bins or groups by descriptor similarity or phenotype modules
      \item Maps higher-order structure from glyph constellations
    \end{itemize}
\end{itemize}

\subsection*{J. Siddhartha Projection and Surface Tracing}
\begin{itemize}[leftmargin=1.5em]
  \item \textbf{Multidimensional Trace:} slices a high-dimensional subfield into 3D using its spread axis.
  \item \textbf{Afterglow:} trailing frames visualized with alpha decay and color gradient.
  \item \textbf{Surface Derivatives:} vectors from grid to surface points used to estimate curvature.
\end{itemize}

\subsection*{K. Surface Extraction — Eitri}
\begin{itemize}[leftmargin=1.5em]
  \item \textbf{Cylindrical Shells:} projected around Brokk axis; grid-aligned.
  \item \textbf{Clock-hand Analysis:} select max-distance vectors per angular bin within tube slices.
  \item \textbf{Multi-Rail Correction:} check across eigenvector axes; update surface candidates if stronger per sector.
\end{itemize}

\subsection*{L. Infrastructure and Instrumentation — F42}
\begin{itemize}[leftmargin=1.5em]
  \item \textbf{F42 Data Model:} flat-typed arrays, no nested allocatables, portable `.txdata` binary format.
  \item \textbf{Patch System:} timestamped changes with semantic messages; Git- and DVC-compatible.
  \item \textbf{Energy Awareness:} Fortran for numeric kernels, no model training, low-carbon compute philosophy.
\end{itemize}

\subsection*{M. Visualization and Documentation — Flyer and Labbook}
\begin{itemize}[leftmargin=1.5em]
  \item \textbf{TOX Flyer:}
    \begin{itemize}
      \item Interactive D3-based visualization
      \item Movie mode: state recording and playback
      \item JavaScript-based scripting and exportable snippets
    \end{itemize}
  \item \textbf{Labbook:}
    \begin{itemize}
      \item Rendered as HTML
      \item Embeds Flyer visualizations or animations
      \item Optional slide deck export (HTML per slide)
    \end{itemize}
\end{itemize}

\end{document}
